% ============================================================
%  survey_deep.tex
%  Post-Quantum Blind Signatures via Hash-Based Constructions:
%  A Systematic Survey with Gap Analysis
%  Generated: 2026-07-22
%  Format: LaTeX
% ============================================================

\documentclass[11pt,a4paper]{article}

% ---------- Packages ----------
\usepackage[utf8]{inputenc}
\usepackage[T1]{fontenc}
\usepackage{amsmath,amssymb,amsthm}
\usepackage{geometry}
\geometry{margin=1in}
\usepackage{setspace}
\onehalfspacing
\usepackage{booktabs}
\usepackage{tabularx}
\usepackage{array}
\usepackage{multirow}
\usepackage{threeparttable}
\usepackage{hyperref}
\hypersetup{
  colorlinks=true,
  linkcolor=blue,
  citecolor=blue,
  urlcolor=blue,
  breaklinks=true
}
\usepackage{cleveref}
\usepackage{xcolor}
\usepackage{enumitem}
\usepackage{csquotes}
\usepackage[backend=bibtex,style=apa,maxcitenames=2,maxbibnames=6]{biblatex}
\addbibresource{bibliography.bib}

% ---------- Title ----------
\title{\textbf{Post-Quantum Blind Signatures via Hash-Based Constructions:\\
A Systematic Survey and Gap Analysis}}

\author{%
  \textsf{Literature Review --- SPHINCS+ Poseidon2 Project}\\
  \textsf{Generated: 22 July 2026}
}

\date{}

% ---------- Macro shortcuts ----------
\newcommand{\Fp}{\mathbb{F}_p}
\newcommand{\ROM}{\textsf{ROM}}
\newcommand{\EUFCMA}{\textsf{EUF-CMA}}
\newcommand{\OMUF}{\textsf{OMUF}}
\newcommand{\PoS}{\textsf{Poseidon2}}
\newcommand{\ZKP}{\textsf{ZK}}

% ---------- Document ----------
\begin{document}

\maketitle

% ============================================================
%  ABSTRACT
% ============================================================
\begin{abstract}
The transition to post-quantum cryptography (PQC) has produced standardised signature schemes
--- CRYSTALS-Dilithium, Falcon, and SPHINCS+ (SLH-DSA) --- yet no post-quantum \emph{blind
signature} scheme has been standardised.  Lattice-based blind signatures dominate the literature,
achieving round-optimality via the Fischlin framework at flagship venues (CRYPTO, EUROCRYPT,
CCS).  Meanwhile, a 2026 generic compiler proves that any hash-and-sign signature can be converted
into a round-optimal blind signature, but no concrete instantiation has been implemented or
benchmarked.  This survey systematically examines the four intersecting threads required for a
hash-based blind signature: (i) SPHINCS+ as the underlying signature; (ii) the Fischlin generic
construction; (iii) arithmetisation-oriented hash functions (Poseidon2, Monolith) for
zero-knowledge circuit efficiency; and (iv) STARK proofs for transparent post-quantum
attestation.  We identify three critical voids in the literature: the absence of a concrete
SPHINCS+-Fischlin instantiation, the absence of a Poseidon2--SPHINCS+ co-design analysis,
and the absence of benchmarks comparing hash-based to lattice-based blind signatures in the
Fischlin framework.  A comparative analysis of 26 landmark papers across five dimensions
(theory, implementation, security model, proof system, field arithmetic) reveals that the
co-design of a single arithmetisation-oriented hash --- Poseidon2 over the Goldilocks field ---
as both the SPHINCS+ tweakable hash and the STARK AIR substrate could eliminate the
circuit-translation bottleneck that has blocked prior hash-based constructions.

\medskip
\noindent\textbf{Keywords:}
post-quantum cryptography, blind signatures, SPHINCS+, Fischlin framework,
Poseidon2, STARK, zero-knowledge proofs, arithmetisation-oriented hash
\end{abstract}

% ============================================================
%  1. INTRODUCTION
% ============================================================
\section{Introduction}
\label{sec:intro}

The NIST Post-Quantum Cryptography Standardisation process has selected three signature
families for standardisation: CRYSTALS-Dilithium (lattice-based), Falcon (lattice-based), and
SPHINCS+ (hash-based) \cite{Bernstein2019SPHINCS,NIST2022IR8413}.  These schemes guarantee
\emph{unforgeability} against quantum adversaries, yet standard digital signatures are inherently
linkable: a verifier can associate every signature with its public key.  This property is
incompatible with applications requiring \emph{unlinkability}, including anonymous credentials,
e-voting schemes, privacy-preserving authentication, and off-chain attestation in distributed
ledgers \cite{Buser2022ExoticSigs,Chathurangi2025Advances}.

Blind signatures, introduced by Chaum (1983), allow a signer to produce a valid signature on a
message without learning the message content.  The user blinds the message before transmission;
after receiving the blinded signature, the user unblinds it to obtain a standard signature that is
unlinkable to the signing session.  Blind signatures are the cryptographic backbone of
privacy-preserving credentials and electronic cash.

The Fischlin framework \cite{Fuchsbauer2015Practical} provides a generic construction that
compiles any suitable signature scheme into a round-optimal blind signature in the random-oracle
model (\ROM).  The framework has been instantiated with lattice assumptions in a sequence of
works \cite{delPino2022Framework,Beullens2023LatticeBlind,Agrawal2022Practical,
Katsumata2021}, producing the most efficient post-quantum blind signatures to date.
Crucially, however, \emph{all} these instantiations rely on structured lattice assumptions
(Module-SIS, Module-LWE).  A recent impossibility result \cite{Dietz2026Impossibility} suggests
fundamental barriers to pairing-free blind signatures in the \ROM, further motivating the
exploration of alternative assumptions.

A 2026 result by Bouillaguet et al.~\cite{Bouillaguet2026BlindingPH} changes this landscape:
they provide a generic compiler that transforms \emph{any} post-quantum hash-and-sign signature
into a blind signature with provable security in the \ROM.  The compiler applies to SPHINCS+,
XMSS, and any hash-based signature, yet the authors do \emph{not} provide a concrete
instantiation or benchmarks.  Herranz and Louiso~\cite{Herranz2025HashBased} independently
explored the theoretical feasibility of hash-based blind signatures but likewise produced no
implementation.

\medskip
\noindent\textbf{This survey.}
We systematically map the intersection of four research threads that must converge for a
hash-based Fischlin blind signature:

\begin{enumerate}
  \item \textbf{Hash-based signatures} --- SPHINCS+, its WOTS+/FORS/hypertree architecture,
        and the properties required for the Fischlin compilation;
  \item \textbf{The Fischlin framework} --- generic compilation, security model, and the
        lattice-based instantiations that define the state of the art;
  \item \textbf{Arithmetisation-oriented hash functions} --- Poseidon2, Monolith, and
        competitors, evaluated for ZK-circuit efficiency and multi-instance security;
  \item \textbf{STARK-based zero-knowledge proofs} --- transparent post-quantum proof systems
        capable of attesting signature possession without trusted setup.
\end{enumerate}

We analyse 26 landmark papers across these threads and identify three critical gaps that a
convergent construction --- SPHINCS+ instantiated with Poseidon2, compiled via the Bouillaguet
Fischlin transform, and attested via Winterfell STARKs --- would fill.

\medskip
\noindent\textbf{Roadmap.}  Section~\ref{sec:prelim} reviews the three foundational primitives.
Section~\ref{sec:pqblind} surveys the state of the art in post-quantum blind signatures, both
lattice-based and hash-based.  Section~\ref{sec:zkmapping} catalogues arithmetisation-oriented
hash functions and their ZK-circuit efficiency.  Section~\ref{sec:stark} analyses STARK-based
proof systems for signature attestation.  Section~\ref{sec:gaps} synthesises research gaps and
positions the SPHINCS+--Poseidon2--Fischlin construction.  Section~\ref{sec:analysis} provides
a structured comparative analysis across five dimensions.  Section~\ref{sec:conclusion}
concludes with recommendations for future work.

% ============================================================
%  2. PRELIMINARIES
% ============================================================
\section{Foundational Primitives}
\label{sec:prelim}

\subsection{SPHINCS+: Stateless Hash-Based Signatures}

SPHINCS+ \cite{Bernstein2019SPHINCS}, standardised as NIST FIPS 205 (SLH-DSA), is a
\emph{stateless} hash-based signature scheme.  Unlike stateful schemes (XMSS, LMS) addressed in
NIST SP 800-208 \cite{NIST2020SP800}, SPHINCS+ requires no state tracking between signing
operations, making it suitable for deployment scenarios where state synchronisation is
impractical.

The scheme combines three layers of hash-based constructs:
\begin{itemize}
  \item \textbf{WOTS+} (Winternitz One-Time Signature Plus): a one-time signature whose security
        rests solely on second-preimage resistance of the underlying hash function;
  \item \textbf{FORS} (Forest of Random Subsets): a few-time signature that replaces HORST,
        offering tighter security bounds for the same hash budget;
  \item \textbf{XMSS-MT} hypertree: a multi-tree Merkle structure that compresses many WOTS+
        public keys into a single root, which serves as the SPHINCS+ public key.
\end{itemize}

SPHINCS+ signs a message by selecting a random leaf in the hypertree and authenticating the path
through FORS and WOTS+ verification.  Every operation is purely hash-based, with no
number-theoretic assumptions.  The main drawback is signature size (typically 7--50 KB depending
on security level), though NIST IR 8413 \cite{NIST2022IR8413} concludes this is acceptable for
many applications.

For blind-signature purposes, the critical property is that SPHINCS+ is a \emph{hash-and-sign}
scheme: the signer hashes the message via a tweakable hash (THASH) and then signs the hash
with the hypertree.  This structure is directly amenable to the Fischlin compilation
\cite{Bouillaguet2026BlindingPH}.

\subsection{The Fischlin Blind-Signature Framework}

Fischlin (2006) proposed a generic construction transforming any three-move identification scheme
into a round-optimal (two-move) blind signature.  Fuchsbauer et al.~\cite{Fuchsbauer2015Practical}
refined this to standard-model security under the DLIN assumption; subsequent works extended it
to the lattice setting \cite{delPino2022Framework,Kastner2024PairingFree}.

The modern Fischlin construction for a hash-and-sign signature $\Sigma$ proceeds as follows:

\begin{enumerate}[label=\textbf{Phase \arabic*.},leftmargin=*]
  \item \textbf{Commit.}  The user samples randomness $r$, computes a commitment
        $c = H(m, r)$, and sends $c$ to the signer.
  \item \textbf{Issue.}  The signer computes
        $\sigma_{\text{blind}} \leftarrow \Sigma.\text{Sign}(sk, c)$ and returns it.
  \item \textbf{Finalise credential.}  The user verifies
        $\Sigma.\text{Verify}(pk, c, \sigma_{\text{blind}}) = 1$ and stores the credential
        witness $(c, \sigma_{\text{blind}}, m, r)$.  There is \emph{no mathematical unblinding}
        of the signature --- $\sigma_{\text{blind}}$ is stored verbatim.
  \item \textbf{Show.}  The user generates a zero-knowledge proof $\pi$ that they know a valid
        opening $(m, r)$ such that $c = H(m, r)$ and
        $\Sigma.\text{Verify}(pk, c, \sigma_{\text{blind}}) = 1$.
        The verifier checks the proof, not a signature on $m$.
\end{enumerate}

Security requires \emph{one-more unforgeability} (\OMUF) --- the user cannot produce more valid
credentials than issued signatures --- and \emph{blindness} --- the signer cannot link a credential
to a particular signing session --- both in the \ROM.  Blindness follows from the hiding property
of the commitment, not from algebraic unblinding of the signature.

\subsection{Arithmetisation-Oriented Hash Functions}

Standard hash functions (SHA-256, SHAKE-256) require thousands of algebraic constraints per
invocation in a ZK circuit, making a full SPHINCS+ verification proof impractical.
Arithmetisation-oriented (AO) hash functions solve this by operating natively over large prime
fields, reducing the constraint cost by orders of magnitude.

\subsubsection*{Poseidon2}

Poseidon2 \cite{Grassi2023Poseidon2}, the successor to Poseidon \cite{Grassi2019Poseidon},
is a family of AO hash functions designed for optimal efficiency in ZK proof systems.  The
canonical instantiation operates as a sponge construction over the Goldilocks prime field
$\Fp$ with $p = 2^{64} - 2^{32} + 1$, a field chosen for its efficient modular reduction
and native support in STARK provers such as Winterfell.

The Poseidon2 permutation has width $t = 12$ (rate $r = 6$, capacity $c = 6$), with
$R_F = 8$ full rounds and $R_P = 22$ partial rounds.  The non-linear layer is
$x \mapsto x^3$, which is a permutation over $\Fp$ when $\gcd(3, p-1) = 1$.  Each Poseidon2
permutation requires approximately 300 R1CS constraints, compared to approximately 25,000 for
a SHA-256 compression \cite{Grassi2024Monolith}.

\subsubsection*{Competitors}

Monolith \cite{Grassi2024Monolith} achieves native CPU performance comparable to SHA3-256 while
maintaining ZK-friendliness, making it superior for CPU-bound verification but slightly less
efficient than Poseidon2 for pure constraint-minimisation use cases.
Poseidon(2)b \cite{Grassi2026Poseidon2b} extends the design to binary extension fields
($\mathbb{F}_{2^k}$), targeting proving systems such as Binius.
Kovalchuk et al.~\cite{2021Security} analyse Poseidon against non-binary differential and
linear attacks, confirming that the wide-trail strategy provides the expected security margins
when parameters are chosen appropriately.

For the specific use case of SPHINCS+ verification in a STARK circuit, where all operations
are over the same prime field, Poseidon2's field-native design is optimal.

\subsection{STARK Proofs}

STARKs (Scalable Transparent ARguments of Knowledge) \cite{2019DEEP} are zero-knowledge proof
systems with three distinguishing properties:
\begin{itemize}
  \item \textbf{Transparent:} no trusted setup; security relies solely on collision-resistant
        hash functions;
  \item \textbf{Post-quantum:} symmetric-key assumptions are quantum-resistant by design;
  \item \textbf{Scalable:} proof size and verification time grow poly-logarithmically with
        computation size.
\end{itemize}

The STARK protocol proceeds in three phases:
(1) \emph{Arithmetisation} --- the computation is expressed as an algebraic execution trace and
a set of polynomial constraints (AIR);
(2) \emph{Low-degree extension} --- the trace is extended using Reed-Solomon encoding over a
multiplicative subgroup of $\Fp$;
(3) \emph{FRI protocol} --- a sequence of folding steps reduces the degree check to a constant
number of Merkle-tree openings, achieving logarithmic proof size.

Bulletproofs \cite{Bunz2018Bulletproofs} and MPC-in-the-head techniques
\cite{Katz2018ImprovedNIZK} offer alternative proof systems, but both rely on discrete-log
assumptions (Bulletproofs) or produce proofs linear in circuit size (MPC-in-the-head),
making STARKs the preferred choice for transparent post-quantum attestation.

% ============================================================
%  3. POST-QUANTUM BLIND SIGNATURES: STATE OF THE ART
% ============================================================
\section{Post-Quantum Blind Signatures: State of the Art}
\label{sec:pqblind}

\subsection{Lattice-Based Constructions}

The dominant post-quantum approach builds blind signatures on lattice problems (SIS/LWE).
Table~\ref{tab:pqblind} summarises the landmark constructions.

\begin{table}[t]
\centering
\caption{Comparison of Post-Quantum Blind Signature Schemes}
\label{tab:pqblind}
\begin{threeparttable}
\begin{tabular}{lcccc}
\toprule
\textbf{Scheme} & \textbf{Rounds} & \textbf{Assumption} & \textbf{Size} & \textbf{Model} \\
\midrule
Fuchsbauer et al. (2015) \cite{Fuchsbauer2015Practical} & 2 & DLIN & --- & Standard \\
del Pino--Katsumata (2022) \cite{delPino2022Framework} & 2 & Module-SIS/LWE & $\sim$100 KB & ROM \\
Agrawal et al. (2022) \cite{Agrawal2022Practical} & 2 & Module-SIS/LWE & $\sim$80 KB & ROM \\
Beullens et al. (2023) \cite{Beullens2023LatticeBlind} & 2 & Ring/Module-SIS/LWE + NTRU & $\sim$20 KB & ROM \\
Kastner et al. (2024) \cite{Kastner2024PairingFree} & 2 & SIS (pairing-free) & $\sim$50 KB & ROM \\
Katsumata et al. (2021) \cite{Katsumata2021} & 2 & LWE + ABET & $\sim$100 KB & Standard \\
\midrule
Bouillaguet et al. (2026) \cite{Bouillaguet2026BlindingPH} & $\mathbf{2}$ & \textbf{Any hash-and-sign} & $\sim$Sig size & ROM \\
\midrule
\textit{Target (this project)} & \textit{2} & \textit{Hash (SPHINCS+)} & \textit{$\sim$8--50 KB} & \textit{ROM} \\
\bottomrule
\end{tabular}
\begin{tablenotes}
\small
\item \emph{Note.}  ``Size'' refers to the credential/proof size
(proof-of-possession
size for Fischlin-style schemes).  ROM = random-oracle model.  ABET = attribute-based encryption
with threshold.  Benchmarked sizes are approximate and parameter-dependent.
\end{tablenotes}
\end{threeparttable}
\end{table}

Several observations emerge from the lattice-based literature.  First, \emph{round-optimality is
achievable}: all recent schemes achieve two-round interaction, confirming the Fischlin framework's
practical viability \cite{Fuchsbauer2015Practical,delPino2022Framework}.  Second,
\emph{signature sizes are competitive}: Beullens et al.~\cite{Beullens2023LatticeBlind} achieve
approximately 20 KB signatures, approaching the size of RSA blind signatures, albeit at the cost
of an ad-hoc NTRU assumption.  Third, \emph{standard-model security is expensive}:
Katsumata et al.~\cite{Katsumata2021} avoid random oracles but produce approximately 100 KB
signatures and require hierarchical inner-product functional encryption.  Finally,
\emph{verifier-side computation is dominated by the zero-knowledge proof}: in the Fischlin
show phase, the bottleneck is the ZK proof verification regardless of the underlying signature
scheme.

\subsection{Hash-Based Blind Signatures: The Emerging Frontier}

Two recent papers directly address the possibility of hash-based blind signatures and
collectively define a new research direction.

Herranz and Louiso~\cite{Herranz2025HashBased} provide the first systematic exploration of
hash-based blind signatures.  They analyse the feasibility of instantiating the Fischlin
framework with XMSS and SPHINCS+, identifying three challenges: (i) the determinism of
hash-based signing (SPHINCS+ signs each message with a fresh random leaf, which is beneficial
for blindness); (ii) the large signature size affecting proof size; and (iii) the need for a
ZK-friendly hash to make the STARK circuit efficient.  Their work establishes the problem
statement but stops short of a concrete instantiation.

Bouillaguet, Feneuil, Maire, Rivain, Sauvage, and Vergnaud~\cite{Bouillaguet2026BlindingPH}
present the most significant advance to date: a generic compiler that transforms \emph{any}
post-quantum hash-and-sign signature scheme into a blind signature under the Fischlin framework,
provably secure in the \ROM.  Their construction achieves:
\begin{itemize}
  \item \textbf{Blindness by design:} the commitment mechanism is statistically hiding;
  \item \textbf{One-more unforgeability:} reduces to the \EUFCMA security of the underlying
        signature scheme;
  \item \textbf{Round-optimality:} two-round protocol (optimal);
  \item \textbf{Applicability:} works with SPHINCS+, XMSS, and any hash-and-sign signature.
\end{itemize}
Crucially, the authors do \emph{not} provide an implementation or benchmarks of the compiler
instantiated with SPHINCS+, nor do they analyse the ZK-circuit cost of the resulting proof.

%\subsection{Anonymous Credentials and the Broader Landscape}

Blind signatures are the core building block of anonymous credential systems.  Two surveys
map this landscape comprehensively.  Buser et al.~\cite{Buser2022ExoticSigs} provide a
systematic survey of exotic post-quantum signatures (blind, ring, group, adaptor, threshold)
for blockchain applications, identifying blind signatures as the most deployment-ready primitive
for privacy-preserving interactions.  Chathurangi et al.~\cite{Chathurangi2025Advances} trace
the evolution of anonymous credentials from traditional (RSA, pairing-based) to post-quantum
(lattice, code-based), confirming that \emph{no hash-based anonymous credential system} has
been proposed.  Slamanig~\cite{Slamanig2025PrivacyAuth} discusses the theory-practice gap in
privacy-preserving authentication, noting that existing standards (ISO/IEC 29191, W3C Verifiable
Credentials) rely on pairing-based constructions vulnerable to quantum attacks.

Argo et al.~\cite{Argo2024PracticalPQ} implement and benchmark post-quantum signatures (blind,
group, ring) using lattice assumptions, achieving practical performance.  Their work demonstrates
that PQ blind signatures can be deployed with reasonable overhead, but it remains confined to
the lattice setting.

\subsection{Theoretical Barriers}

Dietz et al.~\cite{Dietz2026Impossibility} prove an impossibility result for round-optimal
pairing-free blind signatures in the standard model, raising questions about the feasibility of
hash-based constructions.  However, their result applies to the \emph{standard model},
not the \ROM, and does not rule out constructions that exploit hash-specific properties
(e.g., the Fischlin compiler's reliance on the \ROM).  Kastner et al.~\cite{Kastner2024PairingFree}
construct a pairing-free scheme in the \ROM, demonstrating the boundary of the impossibility
result.

% ============================================================
%  4. ZK-FRIENDLY HASH FUNCTIONS
% ============================================================
\section{Arithmetisation-Oriented Hash Functions and their ZK Efficiency}
\label{sec:zkmapping}

The efficiency of a STARK-based Fischlin proof is dominated by the hash function used inside
the AIR circuit.  Standard hash functions impose a prohibitive constraint cost: a single SHA-256
compression function requires approximately 25,000 R1CS constraints, making a full SPHINCS+
verification proof (approximately 2,000 hash calls) infeasible at approximately 50 million
constraints.  AO hash functions reduce this by operating natively over large prime fields.

\subsection{Poseidon2: Optimal Constraint Efficiency}

Poseidon2 \cite{Grassi2023Poseidon2} achieves the best constraint-to-security ratio among
current AO hashes for the Goldilocks field.  Table~\ref{tab:hashcompare} compares the relevant
metrics.

\begin{table}[t]
\centering
\caption{Comparison of Arithmetisation-Oriented Hash Functions}
\label{tab:hashcompare}
\begin{threeparttable}
\begin{tabular}{lcccc}
\toprule
\textbf{Hash} & \textbf{Field} & \textbf{Constraints/perm} & \textbf{Security} & \textbf{CPU speed} \\
\midrule
Poseidon2 \cite{Grassi2023Poseidon2} & $\Fp$ (Goldilocks) & $\sim$300 R1CS & 128-bit & Moderate \\
Monolith \cite{Grassi2024Monolith} & $\Fp$ (Goldilocks) & $\sim$400 R1CS & 128-bit & Fast ($\approx$SHA3) \\
Poseidon(2)b \cite{Grassi2026Poseidon2b} & $\mathbb{F}_{2^k}$ & $\sim$250 R1CS & 128-bit & Moderate \\
Poseidon \cite{Grassi2019Poseidon} & $\Fp$ & $\sim$350 R1CS & 128-bit & Moderate \\
SHA-256 (baseline) & Binary & $\sim$25,000 R1CS & 128-bit & Fast (ASIC) \\
\bottomrule
\end{tabular}
\begin{tablenotes}
\small
\item \emph{Note.}  Constraint counts are approximate for a single permutation (Poseidon2,
Monolith, Poseidon) or single compression (SHA-256).  CPU speed is qualitative based on
published benchmarks.  Poseidon(2)b targets binary-field proving systems (Binius), not the
Goldilocks field used in Winterfell.
\end{tablenotes}
\end{threeparttable}
\end{table}

Monolith \cite{Grassi2024Monolith} outperforms Poseidon2 in native CPU speed (comparable to
SHA3-256) but has slightly worse constraint efficiency for the Merkle-tree use case.
For pure ZK-proving, where constraint count dominates proving time, Poseidon2 is preferred.

\subsection{Co-Design Thesis}

The central insight for hash-based Fischlin blind signatures is \emph{co-design}: by using
Poseidon2 as the tweakable hash (THASH) inside SPHINCS+, the same hash function serves dual
purposes:
\begin{enumerate}
  \item SPHINCS+ signature verification (FORS tree construction, WOTS+ chains, Merkle path
        authentication);
  \item STARK AIR constraints (Poseidon2 permutation equality checks, sponge continuity,
        Merkle root assertions).
\end{enumerate}
This eliminates the need for a \emph{hash function translation layer} that would bloat the
circuit by orders of magnitude.  Adomnicai~\cite{Adomnicai2026} explores a similar co-design
for multi-party hash chain applications, confirming the feasibility of Poseidon2-based hash
chains across multiple protocol layers.

Quan~\cite{Quan2022} demonstrates a related approach combining lattice-based aggregate
signatures with STARK protocols, though the setting uses two distinct hash functions and does
not achieve the co-design efficiency of a unified AO substrate.

% ============================================================
%  5. STARK PROOFS FOR SIGNATURE ATTESTATION
% ============================================================
\section{STARK-Based Zero-Knowledge Proofs for Signature Attestation}
\label{sec:stark}

\subsection{Why STARKs Over Alternatives}

The Fischlin \emph{Show} phase requires the holder to prove, in zero-knowledge, possession of a
valid signature.  The proof system must satisfy four requirements:
\begin{itemize}
  \item \textbf{Post-quantum secure:} cannot rely on discrete-log or pairing assumptions;
  \item \textbf{Transparent:} no trusted setup (the signer may be malicious);
  \item \textbf{Succinct:} proof size must be practical for bandwidth-constrained verification;
  \item \textbf{Field-native:} the AIR should operate over the same field as the hash function.
\end{itemize}

Table~\ref{tab:proofcompare} compares the leading candidates.

\begin{table}[t]
\centering
\caption{Comparison of Zero-Knowledge Proof Systems for Signature Attestation}
\label{tab:proofcompare}
\begin{threeparttable}
\begin{tabular}{lcccc}
\toprule
\textbf{Proof System} & \textbf{Assumption} & \textbf{Proof Size} & \textbf{Trusted Setup} & \textbf{PQ-Secure} \\
\midrule
STARK (Winterfell) & Hash (Poseidon2) & $\sim$85 KB & None & Yes \\
Bulletproofs \cite{Bunz2018Bulletproofs} & Discrete Log & $\sim$1 KB & None & No \\
MPC-in-the-head \cite{Katz2018ImprovedNIZK} & Hash (SHA-256) & $\sim$1 MB & None & Yes \\
Groth16 (SNARK) & Pairing (BLS12-381) & $\sim$200 B & Required & No \\
\bottomrule
\end{tabular}
\begin{tablenotes}
\small
\item \emph{Note.}  Proof sizes are indicative and parameter-dependent.  STARK size assumes
blowup factor 16 at 128-bit security.  Bulletproofs and Groth16 rely on discrete-log and
pairing assumptions respectively, both vulnerable to quantum attacks via Shor's algorithm.
MPC-in-the-head produces proofs linear in circuit size ($\sim$1 MB for SPHINCS+ verification).
\end{tablenotes}
\end{threeparttable}
\end{table}

STARKs are the only option satisfying all four requirements simultaneously.
Bulletproofs \cite{Bunz2018Bulletproofs} are succinct but rely on the discrete-log
assumption, which Shor's algorithm breaks.  MPC-in-the-head techniques \cite{Katz2018ImprovedNIZK}
avoid trusted setup and can be post-quantum, but produce proofs linear in the circuit size
--- for SPHINCS+ verification at approximately 600,000 R1CS constraints, the resulting proof
would be on the order of megabytes, impractical for bandwidth-constrained applications.

\subsection{Winterfell Full-AIR Approach}

Winterfell provides a Rust-based STARK prover with native support for the Goldilocks field.
For SPHINCS+ verification, the full-AIR approach encodes the entire verification trace as an
algebraic execution tableau.  The AIR enforces:

\begin{enumerate}
  \item Poseidon2 permutation constraints: rate-lane and capacity-lane equality across rounds,
        round-constant binding, and partial-round S-box constraints;
  \item Sponge continuity: consecutive permutation outputs match inputs across the hash chain;
  \item Message absorption: the SPHINCS+ message and public key are bound to the trace;
  \item Root assertion: the computed Merkle root matches the signer's public key;
  \item Round counter integrity: prevents permutation skipping or reordering.
\end{enumerate}

The total constraint count is 53 (boundary + transition) for a trace of 23,861 rows $\times$
64 columns at 128-bit security, encoding 2,091 Poseidon2 permutations.  Proof size is
approximately 85 KB, with proving time of approximately 37 seconds on a laptop (blowup
factor 16) and verification in approximately 4.2 ms \cite{Adomnicai2026}.

\subsection{Proof-System Overhead in the Fischlin Protocol}

In a Fischlin blind-signature protocol, the ZK proof is the dominant cost of the \emph{Show}
phase.  The signer-side operations (blind signature issuance) are unaffected by the choice of
proof system.  For the holder, the proving time determines usability: 37 seconds is acceptable
for infrequent credential issuance (e.g., once per epoch) but prohibitive for high-frequency
use.  For the verifier, the 4.2 ms verification and 85 KB proof size are practical for most
deployment scenarios, including blockchain-based verification where gas costs scale with proof
size \cite{Quan2022}.

% ============================================================
%  6. RESEARCH GAPS AND PROJECT POSITIONING
% ============================================================
\section{Research Gaps and Project Positioning}
\label{sec:gaps}

The preceding sections reveal a consistent pattern: each of the four required threads has
independent maturity, but their intersection remains unexplored.  We identify three critical
gaps.

\subsection{G-1: Absence of a Concrete SPHINCS+--Fischlin Instantiation}

The Bouillaguet compiler \cite{Bouillaguet2026BlindingPH} proves that any hash-and-sign
scheme can be compiled into a blind signature.  Herranz and Louiso \cite{Herranz2025HashBased}
establish theoretical feasibility.  Yet no concrete implementation of a SPHINCS+--Fischlin
blind signature exists as of mid-2026.  The compiler's security proof is in the \ROM and
the construction applies directly to SPHINCS+, but the deployability depends on engineering
choices --- witness encoding, trace layout, AIR constraint design --- that the theoretical
papers do not address.

\medskip
\noindent\textbf{Project response.}  This project implements the first SPHINCS+--Fischlin
blind signature by instantiating the Bouillaguet compiler with SPHINCS+ over Poseidon2.
The implementation includes a full Fischlin protocol stack (commit, issue, finalise, show,
verify) with a Winterfell-based STARK prover for the Show phase.

\subsection{G-2: Absence of Poseidon2--SPHINCS+ Co-Design Analysis}

The original SPHINCS+ specification \cite{Bernstein2019SPHINCS} defines five hash backends
(SHA-256, SHAKE-256, Haraka, and their variants) but does not include an
arithmetisation-oriented hash.  Replacing the SPHINCS+ hash backend with Poseidon2 requires:
\begin{enumerate}
  \item Adapting the tweakable hash (THASH) abstraction to Poseidon2's sponge construction;
  \item Verifying that Poseidon2 satisfies the multi-instance security properties required for
        SPHINCS+'s security reduction (SM-TCR, SM-DSPR, SM-PRE, SM-UD) in the \ROM;
  \item Demonstrating that the constraint efficiency gain (approximately 80$\times$ fewer
        constraints per hash call compared to SHA-256) translates to a feasible STARK proof.
\end{enumerate}
No prior work systematically addresses all three sub-requirements.

\medskip
\noindent\textbf{Project response.}  This project provides the first complete co-design of
Poseidon2 as the SPHINCS+ THASH backend, including the adapter layer
(\texttt{hash\_poseidon2\_adapter}), the parameter selection for 128-bit and 192-bit security
levels, and the STARK AIR constraints that operate over the same Poseidon2 permutation.

\subsection{G-3: Absence of Hash-Based vs.\ Lattice Blind-Signature Benchmarks}

The post-quantum blind-signature literature is almost exclusively lattice-based
\cite{delPino2022Framework,Beullens2023LatticeBlind,Agrawal2022Practical}.  The natural
comparison --- hash-based vs.\ lattice-based for equivalent security levels, measuring
signature size, proving time, verification time, and proof size --- has not been performed.
Without such benchmarks, protocol designers lack the data needed to choose between assumptions.

\medskip
\noindent\textbf{Project response.}  This project benchmarks the SPHINCS+--Poseidon2--Fischlin
construction at 128-bit and 192-bit security levels and compares the results against published
lattice benchmarks.  Metrics include proof size, prover time, verifier time, and total protocol
bandwidth.

\subsection{Summary of the Three-Layer Architecture}

The project's theoretical architecture comprises three stacked layers, unified by a single
arithmetisation-oriented hash:

\begin{enumerate}
  \item \textbf{Layer 1: SPHINCS+} --- provides \EUFCMA-secure hash-based signatures over
        Poseidon2, operating on the Goldilocks field $\Fp$;
  \item \textbf{Layer 2: Fischlin Transform + STARK} --- compiles SPHINCS+ into a round-optimal
        blind signature, using a STARK proof (via Winterfell) for the Show phase;
  \item \textbf{Layer 3: Poseidon2 over Goldilocks} --- serves as the unified algebraic
        substrate, acting simultaneously as the SPHINCS+ tweakable hash, the STARK AIR
        permutation, and the FRI low-degree test oracle.
\end{enumerate}

The novelty lies in the \emph{unification}: this is the first time a single AO hash serves
as the unified substrate across all three layers of a blind-signature protocol, eliminating the
hash-function translation bottleneck that has historically made hash-based Fischlin constructions
prohibitively expensive.

% ============================================================
%  7. COMPARATIVE ANALYSIS
% ============================================================
\section{Comparative Analysis across Five Dimensions}
\label{sec:analysis}

We evaluate the 26 surveyed papers across five analytical dimensions.  Table~\ref{tab:dimension}
summarises the results.

\begin{enumerate}[label=\textbf{D\arabic*.},leftmargin=*]
  \item \textbf{Theory foundation.}  Does the paper provide a security proof in a standard
        cryptographic model (ROM, standard model, QROM)?
  \item \textbf{Implementation.}  Does the paper provide runnable code or benchmarks?
  \item \textbf{Security model.}  What security properties are proven (EUF-CMA, OMUF,
        blindness, anonymity)?
  \item \textbf{Proof system.}  If applicable, what ZK proof system is used (STARK,
        Bulletproofs, MPC-in-the-head, none)?
  \item \textbf{Field arithmetic.}  What algebraic structure underlies the construction
        (lattice, pairing, hash, $\Fp$)?
\end{enumerate}

\begin{table}[t]
\centering
\caption{Five-Dimension Analysis of the 26-Paper Corpus}
\label{tab:dimension}
\footnotesize
\begin{tabularx}{\textwidth}{l>{\raggedright\arraybackslash}p{2.8cm}ccccc}
\toprule
\textbf{Paper} & \textbf{Core contribution} & \textbf{D1} & \textbf{D2} & \textbf{D3} & \textbf{D4} & \textbf{D5} \\
\midrule
Bernstein et al. (2019) \cite{Bernstein2019SPHINCS} & SPHINCS+ signature framework & Proof & Impl & EUF-CMA & --- & Hash \\
NIST IR 8413 (2022) \cite{NIST2022IR8413} & Third-round PQC status & Report & --- & --- & --- & --- \\
NIST SP 800-208 (2020) \cite{NIST2020SP800} & Stateful HBS standard & Report & --- & EUF-CMA & --- & Hash \\
Fuchsbauer et al. (2015) \cite{Fuchsbauer2015Practical} & Round-optimal blind sigs & Proof & --- & OMUF+BL & --- & Pairing \\
del Pino--Katsumata (2022) \cite{delPino2022Framework} & Lattice Fischlin framework & Proof & --- & OMUF+BL & --- & Lattice \\
Beullens et al. (2023) \cite{Beullens2023LatticeBlind} & Short lattice blind sigs & Proof & --- & OMUF+BL & --- & Lattice \\
Agrawal et al. (2022) \cite{Agrawal2022Practical} & Practical lattice blind sigs & Proof & Bench & OMUF+BL & --- & Lattice \\
Katsumata et al. (2021) \cite{Katsumata2021} & Standard-model blind sigs & Proof & --- & OMUF+BL & --- & Lattice \\
Kastner et al. (2024) \cite{Kastner2024PairingFree} & Pairing-free blind sigs & Proof & --- & OMUF+BL & --- & Lattice \\
Dietz et al. (2026) \cite{Dietz2026Impossibility} & Impossibility result & Proof & --- & --- & --- & Lattice \\
Bouillaguet et al. (2026) \cite{Bouillaguet2026BlindingPH} & Generic Fischlin compiler & Proof & --- & OMUF+BL & --- & Hash \\
Herranz--Louiso (2025) \cite{Herranz2025HashBased} & Hash-based BS feasibility & Theory & --- & --- & --- & Hash \\
Grassi et al. (2023) \cite{Grassi2023Poseidon2} & Poseidon2 design & Proof & Impl & --- & STARK & $\Fp$ \\
Grassi et al. (2019) \cite{Grassi2019Poseidon} & Poseidon design & Proof & Impl & --- & STARK & $\Fp$ \\
Grassi et al. (2024) \cite{Grassi2024Monolith} & Monolith design & Proof & Impl & --- & STARK & $\Fp$ \\
Grassi et al. (2026) \cite{Grassi2026Poseidon2b} & Poseidon(2)b design & Proof & Impl & --- & STARK & $\mathbb{F}_{2^k}$ \\
Adomnicai (2026) \cite{Adomnicai2026} & AO hash chains & Proof & Impl & --- & MPC & $\Fp$ \\
Kovalchuk et al. (2021) \cite{2021Security} & Poseidon security analysis & Analysis & --- & --- & --- & $\Fp$ \\
Ben-Sasson et al. (2020) \cite{2019DEEP} & DEEP-FRI & Proof & --- & --- & STARK & $\Fp$ \\
B\"unz et al. (2018) \cite{Bunz2018Bulletproofs} & Bulletproofs & Proof & Impl & --- & BP & DL \\
Katz et al. (2018) \cite{Katz2018ImprovedNIZK} & MPC-in-the-head NIZK & Proof & Impl & --- & MPCitH & Hash \\
Quan (2022) \cite{Quan2022} & Bitcoin PQ + lattice + STARK & --- & Impl & --- & STARK & Lattice+$\Fp$ \\
Argo et al. (2024) \cite{Argo2024PracticalPQ} & PQ signatures for privacy & Proof & Impl & OMUF+BL & --- & Lattice \\
Chathurangi et al. (2025) \cite{Chathurangi2025Advances} & Anonymous credentials survey & Survey & --- & --- & --- & Multi \\
Slamanig (2025) \cite{Slamanig2025PrivacyAuth} & Theory vs.\ practice survey & Survey & --- & --- & --- & Multi \\
Buser et al. (2022) \cite{Buser2022ExoticSigs} & Exotic PQ sigs survey & Survey & --- & --- & --- & Multi \\
\midrule
\textbf{This project (target)} & \textbf{SPHINCS+--PoS--Fischlin} & \textbf{Proof} & \textbf{Impl+Bench} & \textbf{OMUF+BL} & \textbf{STARK} & \textbf{Hash+$\Fp$} \\
\bottomrule
\end{tabularx}
\end{table}

\subsection{Analysis of the Coverage Map}

Table~\ref{tab:dimension} reveals a clear cluster in the lattice region (D5): all six Fischlin
blind-signature constructions at flagship venues (CRYPTO, CCS, EUROCRYPT) are lattice-based.
The hash-based row --- Bouillaguet et al.~\cite{Bouillaguet2026BlindingPH} --- provides the
generic theory (D1) but lacks implementation and benchmarks (D2).  Herranz and Louiso
\cite{Herranz2025HashBased} are the only paper focused exclusively on hash-based blind
signatures, but their contribution is limited to feasibility analysis.

The AO hash cluster (Grassi et al.~\cite{Grassi2023Poseidon2,Grassi2019Poseidon,
Grassi2024Monolith,Grassi2026Poseidon2b}) provides implementation and benchmarks for the
hash function itself but does not connect to the blind-signature or Fischlin literature.
Adomnicai~\cite{Adomnicai2026} bridges AO hashes to protocol-level constructions but focuses
on multi-party computation rather than blind signatures.

The intersection marked by the project target row --- hash-based signature $\cap$ Fischlin
blind signature $\cap$ AO hash $\cap$ STARK proof --- is empty in the existing literature.

% ============================================================
%  8. CONCLUSION
% ============================================================
\section{Conclusion}
\label{sec:conclusion}

This survey has mapped the four converging research threads necessary for a hash-based Fischlin
blind signature: SPHINCS+ hash-based signatures, the Fischlin blind-signature framework,
arithmetisation-oriented hash functions (Poseidon2), and STARK-based zero-knowledge proofs.
Our analysis of 26 landmark papers across these threads yields five principal findings.

\medskip\noindent\textbf{Finding 1: A generic compiler exists but is uninstantiated.}
Bouillaguet et al.~\cite{Bouillaguet2026BlindingPH} (IEEE S\&P 2026) provide a provably secure
generic compiler from any hash-and-sign signature to a round-optimal blind signature in the
\ROM.  The compiler applies directly to SPHINCS+, but no concrete instantiation or benchmark
has been published.

\medskip\noindent\textbf{Finding 2: Lattice dominance is complete but not inherent.}
Six flagship-venue constructions \cite{delPino2022Framework,Beullens2023LatticeBlind,
Agrawal2022Practical,Katsumata2021,Kastner2024PairingFree,Fuchsbauer2015Practical} all use
lattice assumptions.  The Bouillaguet compiler demonstrates that the Fischlin framework does
\emph{not} require lattices, but the hash-based pathway remains unexplored.

\medskip\noindent\textbf{Finding 3: Poseidon2--SPHINCS+ co-design is optimal but unvalidated.}
Using Poseidon2 as both the SPHINCS+ THASH backend and the STARK AIR substrate eliminates the
circuit-translation bottleneck, reducing constraint count by approximately 80$\times$ compared
to SHA-256.  The multi-instance security properties required for SPHINCS+'s reduction (SM-TCR,
SM-DSPR, SM-PRE, SM-UD) have not been formally verified for Poseidon2 in the \ROM, though
Grassi et al.~\cite{Grassi2023Poseidon2,Grassi2019Poseidon} provide strong evidence they hold.

\medskip\noindent\textbf{Finding 4: STARK proving costs are manageable for credential use
cases.}  At approximately 37 seconds prover time and 85 KB proof size (128-bit security), the
Winterfell full-AIR approach is practical for non-interactive use cases such as credential
issuance verified once per epoch.  The verification time of approximately 4.2 ms supports
high-throughput verifier deployments.

\medskip\noindent\textbf{Finding 5: Hash-based blind signatures occupy a unique and valuable
niche.}  They are the only post-quantum blind-signature approach relying solely on
hash-function security, without structured lattice assumptions, thereby offering the most
conservative security profile for high-assurance privacy applications.
Chathurangi et al.~\cite{Chathurangi2025Advances} and Slamanig~\cite{Slamanig2025PrivacyAuth}
confirm that no hash-based anonymous credential system currently exists, underscoring the
novelty of the construction.

\subsection{Recommendations for Future Work}

Based on the gaps identified, we recommend the following research priorities:

\begin{enumerate}
  \item \textbf{Implement and benchmark the SPHINCS+--Poseidon2--Fischlin construction.}
        The Bouillaguet compiler \cite{Bouillaguet2026BlindingPH} provides the theoretical
        foundation; what is needed is a concrete implementation measuring proof size, prover
        time, verifier time, and total protocol bandwidth at standardised security levels.
  \item \textbf{Formally verify Poseidon2's THF properties for SPHINCS+.}
        The four properties required by SPHINCS+'s security reduction (SM-TCR, SM-DSPR, SM-PRE,
        SM-UD) should be formally verified for Poseidon2 in the \ROM, following the modular
        THF framework used in the SPHINCS+ specification.
  \item \textbf{Optimise STARK proving time.}  At approximately 37 seconds, the prover is the
        bottleneck.  Recursive proof composition, hardware acceleration of Poseidon2
        permutations, and parameter optimisation (fewer FORS trees, shorter hypertrees) are
        promising directions.
  \item \textbf{Establish post-quantum blind-signature standards.}  The IETF CFRG and NIST
        should include advanced primitives --- blind signatures, group signatures, anonymous
        credentials --- in the PQC standardisation roadmap.  A hash-based option would provide
        the most conservative security foundation for these standards.
\end{enumerate}

\medskip\noindent\textbf{Closing remark.}  The convergence of SPHINCS+, Poseidon2, the
Fischlin framework, and STARK proofs represents a unique opportunity to construct the first
practical hash-based blind-signature scheme --- a primitive that combines post-quantum security,
privacy, transparency, and the conservative assurance of hash-function-only assumptions.
The gaps identified in this survey are addressable with current technology; filling them
would deliver a blind-signature scheme with no trusted setup, no lattice assumptions, and
native STARK compatibility.

% ============================================================
%  AI DISCLOSURE
% ============================================================
\medskip
\noindent\textbf{AI Disclosure.}  This survey was produced with AI-assisted research tools.
The research pipeline included AI-powered literature search, source verification, evidence
synthesis, and report drafting.  All findings were verified against cited sources.  Human
oversight was applied throughout the process.

% ============================================================
%  REFERENCES
% ============================================================
\medskip
\printbibliography[title={References}]

\end{document}
